\documentclass[10pt,a4paper]{article}

\usepackage[margin=20mm]{geometry}
\usepackage{amsmath,amssymb,amsthm}

\theoremstyle{definition}
\newtheorem*{definition*}{Definition}
\newtheorem*{problem*}{Problem}

\newcommand{\CL}{\mathsf{CL}}
\newcommand{\NC}{\mathsf{NC}}
\newcommand{\TC}{\mathsf{TC}}
\newcommand{\SAC}{\mathsf{SAC}}
\newcommand{\CSPACE}{\mathsf{CSPACE}}

\setlength{\parindent}{0pt}
\setlength{\parskip}{5pt}
\setlength{\emergencystretch}{2em}

\begin{document}

\begin{center}
  {\Large\bfseries NC\textsuperscript{2} versus Catalytic Logspace}
\end{center}

\section{Definitions}

\begin{definition*}[Uniform Boolean circuits]
A Boolean circuit is a directed acyclic graph whose internal gates are
$\wedge$, $\vee$, and $\neg$ gates. Unless stated otherwise, $\wedge$ and
$\vee$ have fan-in two. A family $\{C_n\}_{n\geq1}$ is
\emph{logspace-uniform} if a deterministic logspace transducer, given
$(1^n,i)$, outputs the type of the $i$th gate of $C_n$ and its incident wires.
All circuit families below are logspace-uniform and have polynomial size.
\end{definition*}

\begin{definition*}[$\NC^2$ and depth-bounded $\NC$]
The class $\NC^2$ consists of the languages decided by polynomial-size
Boolean circuit families of depth $O(\log^2n)$. More generally,
$\mathsf{NC\mbox{-}DEPTH}(d(n))$ denotes the languages decided by such
families of depth $O(d(n))$.
\end{definition*}

\begin{definition*}[Catalytic logspace]
A catalytic Turing machine is a deterministic Turing machine with a read-only
input tape, an ordinary read-write work tape, and an additional read-write
catalytic tape. For every input $x$ and every string $a$ initially stored on
the catalytic tape, the machine must halt with the correct answer and restore
the catalytic tape exactly to $a$.

Let $\CSPACE(s(n),c(n))$ denote the class of languages decided using
$O(s(n))$ work-tape cells and $O(c(n))$ catalytic-tape cells. Catalytic
logspace is
\begin{equation}
  \CL:=\CSPACE\bigl(O(\log n),n^{O(1)}\bigr).
\end{equation}
No running-time bound is imposed.
\end{definition*}

\begin{definition*}[$\SAC^2$ and clean register programs]
The class $\SAC^2$ is defined like $\NC^2$ except that $\vee$ gates may have
unbounded fan-in, while $\wedge$ gates retain fan-in two.

A register program over a ring $R$ applies a fixed sequence of reversible
register instructions to input elements and to registers with arbitrary
initial contents $\tau_i$. It \emph{cleanly computes} $f$ if each designated
output register finishes with value $\tau_i+f_i(x)$ and every other register
returns to $\tau_i$. Clean register programs can be simulated by catalytic
machines; their number of registers controls catalytic space, while nested or
recursive input calls control ordinary work space.
\end{definition*}

\section{Best Known Result}

The strongest full circuit-class containment in catalytic logspace stated in
the source is
\begin{equation}
  \TC^1\subseteq\CL.
\end{equation}
For bounded-fan-in circuits, the reachable depth extends slightly beyond
logarithmic depth:
\begin{equation}\label{eq:depth}
  \mathsf{NC\mbox{-}DEPTH}(\log n\log\log n)
  \subseteq
  \CL.
\end{equation}
For the full class $\NC^2$, the best simulation supplied by the source is
\begin{equation}\label{eq:nc2}
  \NC^2
  \subseteq
  \CSPACE\left(
    O\left(\frac{\log^2n}{\log\log n}\right),
    n^{O(1)}
  \right).
\end{equation}
This improves the standard $O(\log^2n)$ work-space simulation by a factor
$\Theta(\log\log n)$ while retaining a polynomial-size catalyst. To reach
$\CL$, the work space in~\eqref{eq:nc2} must still be reduced by a factor
$\Theta(\log n/\log\log n)$. Equivalently, at the exact $O(\log n)$
work-space threshold, the depth in~\eqref{eq:depth} must be increased by the
same asymptotic factor to cover depth $O(\log^2n)$.

For every fixed $\varepsilon>0$, the stronger circuit-class simulation
\begin{equation}\label{eq:sac2}
  \SAC^2
  \subseteq
  \CSPACE\left(
    O\left(\frac{\log^2n}{\log\log n}\right),
    2^{O(\log^{1+\varepsilon}n)}
  \right)
\end{equation}
is known. Here the catalyst is superpolynomial, so~\eqref{eq:sac2} is not a
containment in $\CL$. It nevertheless shows that layer merging and polynomial
representations apply beyond bounded-fan-in circuits.

\section{Problem Formalization}

\begin{problem*}[$\NC^2$ versus $\CL$]
Prove the full containment
\begin{equation}
  (\mathsf{Uniform})-\NC^2\subseteq\CL,
\end{equation}
or, failing that, substantially improve the best known simulation by proving:
  \begin{equation}
    (\mathsf{Uniform})-\NC^2
    \subseteq
    \CSPACE\bigl(s(n),n^{O(1)}\bigr)
  \end{equation}
  for some
  \begin{equation}
    s(n)=o\left(\frac{\log^2n}{\log\log n}\right);
  \end{equation}


\end{problem*}

\end{document}
